\subsection{The actual block-growth estimate}
\label{actual-block-growth}

Fix the original place, normalized tower, top lift and
$\mathrm{InDDplus2}$ hypotheses, as in the preceding two estimates.
Let $F$ be an original compact smooth coefficient with a positive
Gram cutoff.  All forms, word operators and Haar measures below are
the prescribed ones.  Write $N_{i,w}$ and $E_{i,w}$ for the normalized
kernel and diagonal of Section~\ref{actual-kernel-fourier}; in
particular their denominator is the full product
$S_w=\prod_{(\varepsilon,a)\in w}n(a)$, with $n(a)>0$.
The finite parameter set $\mathcal Y$ and integers $N,K\ge0$ are
arbitrary.  Admissibility and permitted extensions have exactly
their original definitions.  No restriction on the placement of
the old word is added.\cite{block-growth-source}

\begin{theorem}\label{block-growth-actual}
Let $b\ge1$, $s\ge2$, $j=b+s-1\le r+1$, and let $h$ be an
original valid head of size $b-1$.  Fix an admissible word $w$ with
$|w|+K\le N$.  Suppose that
\[
 \sup_{t\ne0}E_{j-1,w'}(h\mathbin{\|}(s-1,t))<\infty
\]
for every original permitted extension $w'$ of $w$ by at most $K$
letters.  For every complex lattice
$\Lambda=\mathbb Z\omega_1+\mathbb Z\omega_2$ of rank two and
every $z_0\in\mathbb C$, there is a function $k$ on its original
dual lattice $\Lambda^\vee$ such that
\begin{align}
 \frac{\operatorname{covol}(\Lambda)}\pi
       \sum_{z\in z_0+\Lambda}E_{j,w}(h\mathbin{\|}(s,z))
       &\le\sum_{\sigma\in\Lambda^\vee}|k(\sigma)|,
                                                   \label{block-growth-lattice}\\
 |k(\sigma)|&\le
 \limsup_{\substack{t\to0\\t\ne0}}
 \left|N_{j-1,w}\left(h\mathbin{\|}(s-1,t),
        h\mathbin{\|}\bigl(s-1,(t^{-1}+\sigma)^{-1}\bigr)\right)\right|.
                                                   \label{block-growth-limsup}
\end{align}
The sums and limsup have their original extended nonnegative
interpretation.  The limsup is over the full punctured complex
neighborhood, not a ray or a selected sequence.  This is the
original compact \texttt{BlockGrowthBound} supplier.
\end{theorem}

\paragraph{A common compact model for the actual matrix block.}
We first establish the analytic input, retaining the old word.
Put $p=b-1$ and $n=s-1$.  Below the terminal level let
$\rho=r-j\ge0$.  Apply the original $j$th word operator to $F$
first, retaining its full finite auxiliary-index family.  In the
original head quotient, integrate the fixed head $h$ and all
radical coordinates, but leave its final $s\times s$ matrix $Z$
unintegrated.  Denote the resulting coefficient on $G_\rho$ by
$\Phi(Z)$.  This is precisely the source's $\mathrm{PhiFun}$
with a final free block.

Let $\mathcal K\Subset\mathbb C^\times$ be a compact set of
determinant parameters.  Multiply $\Phi(Z)$ by a smooth determinant
cutoff equal to one on $\det Z\in\det(J_n)\mathcal K$ and
supported in a compact subset of $\mathbb C^\times$.  This does
not change either germ orbit considered below.  The resulting
function is compact smooth in $Z$ with values in compact smooth
coefficients on $G_\rho$.  Crucially, those coefficients have
\emph{one common} compact support and positive Gram cutoff.

Here is the properness argument, beyond a pointwise slice assertion.
In the block notation of Section~\ref{actual-partial-head-transfer},
the full head matrix is recovered from the bottom-left block of
the original group matrix.  Eliminate the fixed head by
\[
 Z_{\rm all}=\begin{pmatrix}H&HB\\ AH&AHB+Z\end{pmatrix},
 \qquad A=(Z_{\rm all})_{21}H^{-1},\quad
 B=H^{-1}(Z_{\rm all})_{12}.
\]
The already proved restricted-root Gaussian elimination bounds
$H,H^{-1},A,B$ on the compact support of $F$.  Its denominators
are only the fixed nonzero block determinants of $h$.
The Schur complement $Z$ is therefore bounded.  Its determinant
cutoff bounds $Z^{-1}$.  The explicit inverse radical formulas
\eqref{partial-head-inverse-coordinates} now bound every remaining
integration variable and the residual group matrix.  Thus the
whole preimage is compact, not just each fibre.  The exact identity
\eqref{partial-head-gram} gives a common positive cutoff, since
$|\det Z_{\rm all}|$ is bounded above.  Differentiation under this
common compact integral proves the assertion for every derivative.
The fixed compact actions commute with these coordinates and
preserve their characters, so the residual equivariance is exactly
the original one.

Every old word is covered: its coefficients have denominators only
in the full nonzero head determinant and its fixed, separated
spectral parameters.  On this chart they are smooth.  The operator
was applied before inserting the determinant cutoff, so no new
derivative of that cutoff enters an original word.  Its finite
matrix of residue coefficients, including all rational word
corrections, remains unchanged.

The old positive residue form is continuous on each such common
compact coefficient space in some finite $C^m$ norm.  Indeed the
established finite-regularity Jacquet pairing bound controls the
base form; applying the fixed finite-order residue operator only
increases $m$.  Its rational Gram denominators are bounded by the
common positive cutoff.  This is the actual pairing and its
adjointness theorem (\JacquetUnfoldingReference), not a norm chosen
to replace it.  Complete in that $C^m$ norm if needed.  Thus
$\Phi$ is a compact smooth Banach-valued matrix function to which
Theorem~\ref{terminal-density-compact-germ} applies.

The same construction at $j=r+1$ uses the already proved actual
\texttt{PhiZero} function, followed by the fixed-head partial
average.  It is compact smooth on $\mathrm{GL}_s(\mathbb C)$.
The Schur-complement calculation gives the same common support.
The original positive Gram cutoff is retained.

\paragraph{The vector germ below the terminal level.}
Positivity of the actual old form gives a Hilbert space and vectors
$v(Q)$, for all heads of size $j$, with
\begin{equation}\label{block-growth-low-gram}
                    N_{j,w}(Q,Q')=\langle v(Q),v(Q')\rangle.
\end{equation}
We use an inner product anti-linear in its first variable.
Since the original coefficient form is linear in its first
variable, passage from coefficients to these Gram vectors is
anti-linear.  In particular it conjugates the germ character.
This also follows by expanding the squared norm of every finite
linear combination using the exact original kernel.

Set $M=k+2\rho+3$ and let $a_h$ be the unchanged head factor in
$n\mathrm{Pt}$.  For the two block patterns its defining product
gives exactly
\begin{align*}
 n\mathrm{Pt}(h\mathbin{\|}(n,t),(1,c))
    &=a_h\bigl(|t|^{n(M+s-2)}|c|^{M-2}\bigr)^2,\\
 n\mathrm{Pt}(h\mathbin{\|}(s,z))
    &=a_h|z|^{2s(M+s-3)}.
\end{align*}
The empty fixed head is included, with $a_h=1$.
Put $\mu_h=\sqrt{a_h}/\sqrt{S_w}$.  At $t=x^{-1}$ and
$c=x^nD$, the first square root is
$\sqrt{a_h}|D|^{M-2}|t|^{ns}$.  The factor
$|x|^n$ therefore changes it into
$\sqrt{a_h}|D|^{M-2}|t|^{n^2}$.

The source planar measure is Lebesgue measure.  For
$\beta_U=2i$, the self-dual measure in the compact matrix germ is
$d\mu_{\beta_U}=\pi^{-1}d\operatorname{Re}z\,
d\operatorname{Im}z$.  The two-block and long-block quotient
dimensions are $d_B=n(n+3)/2$ and $d_L=n(n+1)/2$ respectively.
Thus the two source integrals are $\pi^{d_B}\Omega_B$ and
$\pi^{d_L}\Omega_L$, with $d_B-d_L=n$.  The previously proved
block factorization has Jacobian one and splits the actual
character into the old character and the two new corner terms.
It follows that, uniformly for $D\in\mathcal K$,
\begin{equation}\label{block-growth-low-germ}
 \begin{split}
 \pi^{-n}|x|^n
       v(h\mathbin{\|}(n,x^{-1}),(1,x^nD))
  =\sum_{z^s=(-1)^{s+1}D}
       \frac{\psi_{-\beta_U}(szx)}{s|z|^n}
                  v(h\mathbin{\|}(s,z))+O_{\mathcal K}(|x|^{-1}).
 \end{split}
\end{equation}
This is an equality with error in the actual Gram norm.
For the scalar check, the long-block power after dividing by
$s|z|^n$ is $s(M+s-3)-n=s(M-2)+n^2$, exactly the power in the
matrix germ after $|D|=|z|^s$ is substituted.  Both sides carry
the same $\mu_h$ and the same old residue matrix.

The compact-germ proof is uniform on $D\in\mathcal K$: its
cofactor lower bound, support bounds, distinct critical roots and
nondegenerate Hessians are uniform on such a compact set.
No rank-uniform remainder constant is required here.  For small
$|x|$ the two-block integral vanishes, because its matrix contains
the unbounded entry $x^{-1}$ whereas $\Phi$ has compact support.
For intermediate $|x|$ common compactness gives a bound.
Increasing the constant therefore proves
\eqref{block-growth-low-germ} with an $O(|x|^{-1})$ bound for
\emph{all} $x\ne0$.

\paragraph{The terminal vector germ, including its density.}
At $j=r+1$, the actual terminal type decomposition and positivity
give finitely many nonnegative numbers $c_i$, compact coefficients
$F_i$, and fixed linear maps $L_i$ such that, for
\[
 v_T(Q)=\bigl(\sqrt{c_i}\,\omega(Q)L_iV_{F_i}(Q)\bigr)_i,
 \qquad
 \omega(Q)=\sqrt{n\mathrm{Pt}(Q)}
          \sqrt{\left(\frac{|\det Q|^2}{G(Q)}\right)^{|w|}},
\]
one has
\begin{equation}\label{block-growth-top-gram}
              N_{r+1,w}(Q,Q')=\langle v_T(Q'),v_T(Q)\rangle.
\end{equation}
For $w\ne\varnothing$ the index is exactly $(e,\rho)$, where
$e$ is the original tensor index and $\rho$ is an original compact
type.  The data are
\[
 c_{e,\rho}=\frac{\bigl(\sqrt{\operatorname{Re}c_W(w,e,\rho)}\bigr)^2}{S_w},
 \qquad F_{e,\rho}=\Theta_{w,e}F,\qquad
 L_{e,\rho}=\operatorname{hsVec}\circ(E_\rho\circ\,\cdot\,),
\]
using the compact representative of $\Theta_{w,e}F$ already
constructed.  For the empty word there is one index, $c=1$,
$F_i=F$ and $L_i=\operatorname{hsVec}$.  Thus the formula retains
the actual spectral factors and $S_w$, and is the original
definition of the top kernel.  The all-head terminal identity
of Theorem~\ref{terminal-transfer-every-head} identifies its
singleton fibres with the prescribed preceding integral.
No finite type projection of a different coefficient is used.
Each weighted summand is defined to be zero where its original
vector is zero; its displayed quotients are evaluated only on
the positive Gram support of that vector.

Write $T_*=TT'>0$.  For the two-block point the exact Gram and
word factors give, on its nonzero support,
\[
 \rho(x,D)=
 \left(1-\frac{T_*}{|x^nD|^2}\right)^{(k-1-|w|)/2}.
\]
Indeed $G=|\det h|^2|x^{-1}|^{2n}(|x^nD|^2-T_*)$ and
$|\det Q|^2/G=|x^nD|^2/(|x^nD|^2-T_*)$.
For the long block the latter ratio is one.  The point weights
are
\[
 \begin{split}
 n\mathrm{Pt}(h\mathbin{\|}(n,t),(1,c))
       &=a_h|t|^{2n(k+n)}(|c|^2-T_*)^{k-1},\\
 n\mathrm{Pt}(h\mathbin{\|}(s,z))
       &=a_h|z|^{2s(k+s-2)}.
 \end{split}
\]
These identities are used only where the coefficient is nonzero;
the original positive Gram cutoff makes every denominator there
positive and bounded away from zero on parameter compact sets.
For large $|x|$ and $D\in\mathcal K$,
$\rho=1+O_{\mathcal K}(|x|^{-2n})$.  Apply the compact matrix
germ to each actual \texttt{PhiZero} coefficient, keeping the same
Lebesgue conversions.  The exponent identity
$s(k+s-2)-n=s(k-1)+n^2$ gives
\begin{equation}\label{block-growth-top-germ}
 \begin{split}
 |x|^n v_T(h\mathbin{\|}(n,x^{-1}),(1,x^nD))
   =\sum_{z^s=(-1)^{s+1}D}
       \psi_{\beta_U}(szx)\frac{\pi^n}{s|z|^n}
                  v_T(h\mathbin{\|}(s,z))+O_{\mathcal K}(|x|^{-1}).
 \end{split}
\end{equation}
The extra factor $\rho-1$ contributes $O(|x|^{-2n})$ times
the bounded stationary-phase main expression, hence is absorbed
because $n\ge1$.  For bounded $x$, the actual support cutoff
handles any apparent zero of $G$; for $x\to0$ compact matrix
support again makes the two-block coefficient zero.
Thus the displayed bound also holds for all $x\ne0$.

\paragraph{One Hilbert-space argument with the exact two scalars.}
Here are the common data supplied by the two preceding paragraphs.
There are a Hilbert space $H$, continuous functions
$u_x(D)$ and $b(z)$ away from their zero arguments, and
\[
 \epsilon=(-1)^{s+1},\qquad
 B(z)=s|z|^{s-1}b(z),\qquad
 A_x(D)=\sum_{z^s=\epsilon D}\psi_\beta(szx)b(z),
\]
such that $u_x-A_x=O_{\mathcal K}(|x|^{-1})$ on every
$\mathcal K\Subset\mathbb C^\times$ and
$\sup_{x\ne0}\|u_x\|_{L^2(dD)}^2<\infty$.
Put
\[
 U_x(c)=|x|^{-n}u_x(c/x^n),\qquad
 A^c_x(c)=|x|^{-n}A_x(c/x^n).
\]
Writing $E(z)=E_{j,w}(h\mathbin{\|}(s,z))$ and
$N(t,t')=N_{j-1,w}(h\mathbin{\|}(n,t),h\mathbin{\|}(n,t'))$,
the \emph{exact} relations are
\begin{equation}\label{block-growth-data}
 E(z)=\kappa_2\|B(z)\|^2,\qquad
 |N(t,t')|=\kappa_1|\langle U_{1/t},U_{1/t'}\rangle_{L^2}|,
 \qquad \mathrm{fsv}(\sigma,z)=\psi_\beta(\eta\sigma z),
\end{equation}
Here $Q_x(D)=h\mathbin{\|}(n,x^{-1}),(1,x^nD)$ and
\[
 \begin{aligned}
 j\le r:&\quad u_x(D)=\pi^{-n}|x|^nv(Q_x(D)),
       &&b(z)=\frac{v(h\mathbin{\|}(s,z))}{s|z|^n},\\
 j=r+1:&\quad u_x(D)=|x|^nv_T(Q_x(D)),
       &&b(z)=\frac{\pi^nv_T(h\mathbin{\|}(s,z))}{s|z|^n}.
 \end{aligned}
\]
The signs and two scalars are
\[
\begin{array}{c|c|c}
 & (\beta,\eta)&(\kappa_1,\kappa_2)\\ \hline
 j\le r&(-\beta_U,1)&(\pi^{2n},1)\\
 j=r+1&(\beta_U,-1)&(1,\pi^{-2n}).
\end{array}
\]
The Gram relations in \eqref{block-growth-data} follow by
polarizing the already proved exact head integral, including
its terminal version.  Under $c=x^nD$ the Jacobian is
$|x|^{2n}$, so the displayed dilation is unitary on $L^2$.
The assumed bounded diagonal for the extension $w'=w$ therefore
gives the asserted uniform $L^2$ bound.  The hypotheses on the
other extensions are preserved, though not needed for this step.
In both cases $0\le\kappa_2\le\pi\kappa_1$.

We give the remaining argument explicitly.  On a compact annulus
in $D$, expand $\|A_x(D)\|^2$.  Diagonal root pairs give
$\int\sum_{z^s=\epsilon D}\|b(z)\|^2\,dD$.
Every off-diagonal root pair has the nonconstant linear phase
$\psi_\beta(s(\zeta-1)zx)$, where $\zeta^s=1$, $\zeta\ne1$.
The change of variables $D=\epsilon^{-1}z^s$ and the
Riemann--Lebesgue lemma make those integrals tend to zero as
$|x|\to\infty$.  The compact germ and the uniform $L^2$ bound,
followed by monotone exhaustion of the annuli, imply
\[
 \int_{\mathbb C}\|B(z)\|^2\,dz<\infty.
\]
The Jacobian here is exactly $s^2|z|^{2s-2}$.
Consequently $A_x\in L^2$ for every $x$, with
$\|A_x\|_2^2\le s\|B\|_2^2$.  In particular $E$ is integrable.

For fixed $\delta\in\mathbb C$, one has
\begin{equation}\label{block-growth-packet-limit}
 \begin{split}
 \langle A^c_x,A^c_{x+\delta}\rangle
      &\longrightarrow\widehat B(\delta)
          :=\int\psi_\beta(\delta z)\|B(z)\|^2\,dz,\\
 \langle U_x-A^c_x,A^c_{x+\delta}\rangle&\longrightarrow0,
 \qquad
 \langle A^c_x,U_{x+\delta}-A^c_{x+\delta}\rangle
      \longrightarrow0.
 \end{split}
\end{equation}
To verify the first limit, in the second root packet choose
the branch $q_x=(x/(x+\delta))^{n/s}$ tending to one.  Corresponding
roots differ by $q_x$, and their diagonal phase is
\[
 sz\bigl((x+\delta)q_x-x\bigr)\longrightarrow\delta z.
\]
Distinct roots give a linear phase whose frequency tends to
infinity.  On a compact annulus, uniform convergence of the
amplitudes and Riemann--Lebesgue prove the limit.  The omitted
tails are uniformly small by Cauchy--Schwarz and the $L^2$ tails
of $B$; the dilations $q_x$ stay near one.  The same truncation,
the local germ bound, and the uniform norms of $U_x,A^c_x$ prove
the other two limits.  This proves the assertions over the full
complex $|x|\to\infty$ filter, with no angular restriction.

\paragraph{The positive residual term and the punctured limsup.}
Fix the rank-two dual lattice $\Gamma=\Lambda^\vee$.
Take a translation-invariant mean $\mathfrak m$ on bounded
functions on $\Gamma$ which agrees with limits at infinity.
For clarity, average over larger integer-coordinate squares
and take a weak-star cluster point.  The boundary-to-volume
ratio tends to zero, so the cluster point is translation
invariant; every finite set has mean zero.  Set $U_0=A^c_0=0$
if necessary and $R_x=U_x-A^c_x$.  All these vectors are uniformly
bounded in $L^2$.

The function
\[
 C(\delta)=\mathfrak m_x\langle R_x,R_{x+\delta}\rangle
       \qquad(\delta\in\Gamma)
\]
is positive definite.  Indeed translation invariance turns
each finite quadratic sum into the mean of a squared Hilbert
norm.  Equation~\eqref{block-growth-packet-limit} gives
\[
 \mathfrak m_x\langle U_x,U_{x+\delta}\rangle
                  =\widehat B(\delta)+C(\delta).
\]
Define on $\Gamma$
\begin{equation}\label{block-growth-k}
 A(\sigma)=\frac{\kappa_2}\pi C(\eta\sigma),\qquad
 k(\sigma)=\frac1\pi\int\mathrm{fsv}(\sigma,z)E(z)\,dz
                       +A(\sigma).
\end{equation}
Then $A$ is positive definite and
\[
 |k(\sigma)|\le\frac{\kappa_2}\pi
       \limsup_{\substack{|x|\to\infty\\x\in\Gamma}}
                      |\langle U_x,U_{x+\eta\sigma}\rangle|.
\]
Use $\kappa_2/\pi\le\kappa_1$ and
\eqref{block-growth-data}.  For $\eta=1$, put $t=x^{-1}$.
For $\eta=-1$, interchange the two Hilbert vectors inside
the absolute value and put $t=(x-\sigma)^{-1}$.
In either case the second parameter is precisely
$(t^{-1}+\sigma)^{-1}$.  The finitely many undefined reciprocal
values have mean zero.  The limsup along these lattice parameters
is bounded above by the full punctured complex limsup, proving
\eqref{block-growth-limsup}.  The residual term has not been
discarded, and no pointwise limit of its Gram kernels is assumed.
Extend $A$ and $k$ by zero off $\Gamma$ if functions on all of
$\mathbb C$ are required, as in the original contract.

\paragraph{Poisson summation with the positive residual term.}
The function $E$ is continuous off zero, by the common compact
models and their actual weights.  At zero it is zero, since the
source coefficient of an invalid head is zero.  Positivity therefore
makes $E$ lower semicontinuous everywhere.  On the torus
$\mathbb C/\Lambda$ set
\[
        f(z)=\frac{\operatorname{covol}(\Lambda)}\pi
                         \sum_{\lambda\in\Lambda}E(z+\lambda).
\]
It is nonnegative, integrable and lower semicontinuous.  For
probability Haar measure on that torus its Fourier coefficient is
$\pi^{-1}\int\mathrm{fsv}(\sigma,z)E(z)\,dz$.
These normalizations
explain the exact factor in \eqref{block-growth-lattice}.

Use the product Fejer kernels in any integer basis of $\Gamma$.
Their coefficient weights are
$a_L(m,n)=(1-|m|/(L+1))_+(1-|n|/(L+1))_+$.
Positive definiteness of $A$ gives
\[
       \sum_{\sigma\in\Gamma}a_L(\sigma)A(\sigma)
                           \overline{\mathrm{fsv}(\sigma,z)}\ge0:
\]
the sum is the finite positive-definite quadratic form for
indices in $\{0,\ldots,L\}^2$, divided by $(L+1)^2$.
It follows from \eqref{block-growth-k} that
\[
       (\mathcal F_L*f)(z_0)
        \le\left|\sum_\sigma a_L(\sigma)k(\sigma)
                          \overline{\mathrm{fsv}(\sigma,z_0)}\right|
        \le\sum_\sigma|k(\sigma)|.
\]
If the last sum is infinite the desired inequality is automatic.
Otherwise the same finite-sum inequality applies for every $L$.
Positivity and concentration of the Fejer kernels, together with
lower semicontinuity, give
$f(z_0)\le\liminf_L(\mathcal F_L*f)(z_0)$ at \emph{every} point.
This proves \eqref{block-growth-lattice}, including the rational
sampling points and any lattice translate containing zero.

The construction is for every original admissible old word and
every original head in the theorem.  Its auxiliary compact-germ
constants may depend on these fixed data; they disappear entirely
from the final inequality.  In rank $r=0$ the block range is empty.
This closes the actual compact block-growth supplier.  Extension
to the final finite-decay coefficient, the arithmetic construction
of its top lift and integral lattices, and the complete global
rational-return estimate remain separate obligations.
